%% LyX 2.4.4 created this file.  For more info, see https://www.lyx.org/.
%% Do not edit unless you really know what you are doing.
\documentclass[english]{article}
\usepackage[T1]{fontenc}
\usepackage[latin9]{inputenc}
\usepackage{enumitem}
\usepackage{amsmath}
\usepackage{amsthm}
\usepackage{cancel}
\usepackage{geometry}
\geometry{verbose,tmargin=1in,bmargin=1in,lmargin=1in,rmargin=1in}

\makeatletter
%%%%%%%%%%%%%%%%%%%%%%%%%%%%%% Textclass specific LaTeX commands.
\numberwithin{equation}{section}
\numberwithin{figure}{section}
\newlength{\lyxlabelwidth}      % auxiliary length 

%%%%%%%%%%%%%%%%%%%%%%%%%%%%%% User specified LaTeX commands.
\usepackage{fancyhdr}
\usepackage{lastpage}
\pagestyle{fancy}
\usepackage{enumitem}
\usepackage{ifthen}
%sets math fonts to be more readable
\everymath=\expandafter{\the\everymath\displaystyle}
%\DeclareMathSizes{10}{10}{10}{10}
\usepackage{multicol}

\usepackage{tikz}
\usetikzlibrary{plotmarks}
\usepackage{pgfplots}
\pgfplotsset{compat=newest}

\usepackage{calc}
\usepackage[nomessages]{fp}

\usepackage{advdate}    % Advancing/saving dates
\usepackage{datetime}   % Dates formatting
\usepackage{datenumber} % Counters for dates
% Added by lyx2lyx
\setlength{\parskip}{\smallskipamount}
\setlength{\parindent}{0pt}

\makeatother

\usepackage{babel}
\begin{document}
\renewcommand{\author}{Dr.\,James Ashe}

\newcommand{\classNum}{131}

\newcommand{\classSec}{B}

\newcommand{\nameType}{Name:}

\newcommand{\examType}{Worksheet 3.1}

\renewcommand{\date}{\formatdate{14}{10}{2013}}

%%First page's header%%%%%%%%%%%%%%%%%%%%%%
\fancypagestyle{firststyle} %%header settings for first pagr
{
	\fancyhead[L]{MTH \classNum{}(\classSec{}) \author{}\\
	\textbf{\examType{}} \date{}}
	\fancyhead[C]{\textbf{\nameType}}
	\setlength{\headsep}{14pt}
}
%%%%%%%%%%%%%%%%%%%%%%%%%%%%%%%%%%%%%%%%%%

%%Next page's header%%%%%%%%%%%%%%%%%%%%%%%
\setlist{nolistsep}
\lhead{\classNum{}(\classSec{})} 
\chead{\textbf{\nameType}} 
\cfoot{\thepage{}~of~\pageref{LastPage}} 
\setlength{\headsep}{5pt} 
%%%%%%%%%%%%%%%%%%%%%%%%%%%%%%%%%%%%%%%%%%%

%%Various settings; see the preamble for other settings%%%%%
%%\setenumerate[0]{leftmargin=12pt,itemindent=0pt,labelwidth=0pt}
\renewcommand{\labelenumi}{\textbf{\arabic{enumi}.}}
\renewcommand{\labelenumii}{\textbf{\alph{enumii}.}}
\thispagestyle{firststyle}
%%%%%%%%%%%%%%%%%%%%%%%%%%%%%%%%%%%%%%%%%%
\newcommand{\xmax}{10}
\newcommand{\xmin}{-10}
\newcommand{\xstep}{0} %must be xmin+n
\newcommand{\ymax}{10}
\newcommand{\ymin}{-10}
\newcommand{\ystep}{0} %must be ymin+n

\newcommand{\xspace}{\vspace{1.2cm}}\begin{multicols}{2}
\begin{enumerate}[resume]
\item Write the function's equation from the graph.

\begin{enumerate}[resume]
\item []\renewcommand{\xmax}{10} \renewcommand{\xstep}{5}
\renewcommand{\ymax}{\xmax} \renewcommand{\ystep}{5}
%%%%%%%
\begin{tikzpicture} 
	\begin{axis}[
		axis x line=middle, axis y line=middle,     		
		xtick={-\xmax,-\xstep,...,\xmax},
		ytick={-\ymax,-\ystep,...,\ymax},
		minor x tick num={4},
		minor y tick num={4},
		grid=both,
		xlabel={$$},     
		ylabel={$$},
		xmin=-\xmax.25,xmax=\xmax.25,
		ymin=-\ymax.25,ymax=\ymax.25,     		width=6cm,     		height=6cm] 		\addplot[red,very thick,mark=noner,smooth,domain=\xmin:\xmax]{(x-3)*(x-3)-4}; 
	\end{axis}
\end{tikzpicture}
\end{enumerate}
\end{enumerate}
\end{multicols}\vfill

\begin{multicols}{2}
\begin{enumerate}[resume]
\item [\textbf{2.}]Write the function's equation from the graph.

\begin{enumerate}[resume]
\item []\renewcommand{\xmax}{10} \renewcommand{\xstep}{5}
\renewcommand{\ymax}{\xmax} \renewcommand{\ystep}{5}
%%%%%%%
\begin{tikzpicture} 
	\begin{axis}[
		axis x line=middle, axis y line=middle,     		
		xtick={-\xmax,-\xstep,...,\xmax},
		ytick={-\ymax,-\ystep,...,\ymax},
		minor x tick num={4},
		minor y tick num={4},
		grid=both,
		xlabel={$$},     
		ylabel={$$},
		xmin=-\xmax.25,xmax=\xmax.25,
		ymin=-\ymax.25,ymax=\ymax.25,     		width=6cm,     		height=6cm] 		\addplot[red,very thick,mark=noner,smooth,domain=\xmin:\xmax]{-1*(x+2)*(x+2)+2}; 
	\end{axis}
\end{tikzpicture}
\end{enumerate}
\end{enumerate}
\end{multicols}\vfill

\begin{multicols}{2}
\begin{enumerate}[resume]
\item [\textbf{3.}]Let $f(x)=(x-3)^{2}+1$.

\begin{enumerate}[resume]
\item Graph $f(x)$, clearly label the vertex and the $y-$intercept.\vfill
\item What is the maximum of $f(x)$?\vfill
\item What is the minimum of $f(x)$?\vfill
\item Find the domain and range of $f(x)$.\vfill\columnbreak
\item []\renewcommand{\xmax}{10} \renewcommand{\xstep}{5}
\renewcommand{\ymax}{\xmax} \renewcommand{\ystep}{5}
%%%%%%%
\begin{tikzpicture} 
	\begin{axis}[
		axis x line=middle, axis y line=middle,     		
		xtick={-\xmax,-\xstep,...,\xmax},
		ytick={-\ymax,-\ystep,...,\ymax},
		minor x tick num={4},
		minor y tick num={4},
		grid=both,
		xlabel={$$},     
		ylabel={$$},
		xmin=-\xmax.25,xmax=\xmax.25,
		ymin=-\ymax.25,ymax=\ymax.25,     		width=6cm,     		height=6cm] 		%\addplot[red,very thick,mark=noner,smooth,domain=\xmin:\xmax]{-1*(x+2)*(x+2)+2}; 
	\end{axis}
\end{tikzpicture}
\end{enumerate}
\end{enumerate}
\end{multicols}\vfill
\end{document}
